\documentclass[preview]{standalone}

\usepackage{amsmath}
\usepackage{amssymb}
\usepackage{stellar}
\usepackage{definitions}
\usepackage{bettelini}

\begin{document}

\id{matrix-iterative-methods}
\genpage

\section{Matrix iterative methods}

\begin{snippettheorem}{matrix-power-convergence-condition}{Matrix power convergence}
    Sia \(P \in \matrices_{n\times n}(\complexnumbers)\) e \(k \in \naturalnumbers\). Allora,
    \[
        \lim_k P^k = 0
    \]
    se e solo se \(\rho(P) < 1\).
\end{snippettheorem}

\begin{snippetproof}{matrix-power-convergence-condition-proof}{matrix-power-convergence-condition}{Matrix power convergence}
    \begin{iffproof}
        \item Sappiamo che \[
            \lim_k P^k = 0
        \]
        e consideriamo \(\lambda\) autovalore
        di \(P\) con autovettore associato \(x \neq 0\).
        \begin{enumerate}
            \item Da un lato sappiamo che \(P^k x = \lambda^k x\);
            \item Dall'altro sappiamo che 
            \[
                \lim_k (P^k y) = \left(
                    \lim_k P^k
                \right) y = 0, \quad \forall y \in \complexnumbers^n
            \]
        \end{enumerate}
        Ciò implica che
        \[
            \lim_k \lambda^k x = 0
        \]
        che necessita \(|\lambda| < 1\).
        \item Consideriamo separatamente il caso diagonalizzabile e quello non diagonalizzabile:
        \begin{enumerate}
            \item \emph{Caso diagonalizzabile:}
            Abbiamo \(P = TDT^{-1}\)
            dove \(D\) è la matrice diagonale degli autovalori \(\lambda_i\).
            Elevando alla \(k\) abbiamo quindi \[
                P^k = (TDT^{-1})^k = T D^k T^{-1}
            \]
            Quindi, per continuità del prodotto matriciale, abbiamo
            \[
                \lim_k P^k = T \left(\lim_k D^k\right) T^{-1}
            \]
            Ma dalla forma di \(D\), il limite tende a \(0\)
            se e solo se \(|\lambda_i| < 1\).
            Ma siccome il raggio spettrale è minore di \(1\) allora vale.
            \item \emph{Caso non diagonalizzabile:}
            Utilizzando la scomposizione di Schur otteniamo
            \(P = QTQ^*\) con \(Q\) unitaria e \(T\) triangolare superiore.
            Abbiamo
            \[
                \lim_k P^k = 0
            \]
            se e soltanto se
            \[
                \lim_k T^k = 0
            \]
            in quanto \(P^k = QT^k Q^*\).
            Notiamo che \[
                0_{\matrices_{n \times n}(\complexnumbers)} \leq |T|
                = \begin{pmatrix}
                    |\lambda_1| &  &  & \ast \\
                    & |\lambda_2| & &  \\
                    &  & \ddots &  \\
                    0 &  &  & |\lambda_n| \\
                \end{pmatrix} \leq 
                \begin{pmatrix}
                    \gamma_1 &  &  & \ast \\
                    & \gamma_2 &  &  \\
                    &  & \ddots &  \\
                    0 &  &  & \gamma_n \\
                \end{pmatrix} = \hat T
            \]
            Stiamo usando l'ordinamento parziale indotto sulle matrici.
            Nella forma di Schur, possiamo permutare gli elementi della diagonale come vogliamo,
            quindi scegliamo di ordinarli in modo crescente:
            \[
                |\lambda_1| \leq |\lambda_2| \leq \dots \leq |\lambda_n|
            \]
            L'idea è la seguente: affinché \(\hat T\) sia sicuramente diagonalizzabile, vogliamo che i suoi elementi diagonali \(\gamma_i\) siano a due a due distinti. 
            Ogni qualvolta troviamo dei termini in \(|T|\) che coincidono, aggiungiamo un valore ad uno tale da non fargli superare il successivo nella catena.
            Ad esempio, se abbiamo un blocco di autovalori con lo stesso modulo:
            \[
                |\lambda_j| = |\lambda_{j+1}| = \cdots = |\lambda_q| < |\lambda_{q+1}|
            \]
            poniamo \(\varepsilon = |\lambda_{q+1}| - |\lambda_j|\) e definiamo i nuovi valori aggiungendo incrementi successivi di \(\frac{\varepsilon}{N}\) (con \(N\) opportunamente grande):
            \[
                \gamma_j = |\lambda_j|, \quad \gamma_{j+1} = |\lambda_j| + \frac{\varepsilon}{N}, \quad \gamma_{j+2} = |\lambda_j| + \frac{2\varepsilon}{N}, \dots
            \]
            facendo attenzione a restare strettamente minori del valore distinto successivo.
            Imponiamo infine che l'ultimo termine \(\gamma_n < 1\) (il che è possibile poiché \(\rho(P) < 1\)), in maniera tale da avere tutti numeri non negativi, a due a due distinti, e tutti minori di \(1\).\\
            Notiamo che \(|T^k| \leq |T|^k\) per espansione del prodotto e applicando la disuguaglianza triangolare.
            Passando all'esponente per \(k\), abbiamo la seguente catena:
            \[
                0_{\matrices_{n \times n}(\complexnumbers)} \leq |T^k| \leq |T|^k \leq \hat T^k
            \]
            Ma siccome \(\hat T\) ha autovalori tutti distinti è diagonalizzabile, e avendo \(\rho(\hat T) < 1\), per il caso precedente il suo limite tende a zero. 
            Per il teorema del confronto, il limite della catena collassa a \(0\), implicando \(\lim_k |T^k| = 0\), e dunque \(\lim_k P^k = 0\).
        \end{enumerate}
    \end{iffproof}
\end{snippetproof}

\end{document}